%%%%%%%%%%%%%%%%%%%%%%%%%%%%%%%%%%%%%%%%%
% Short Sectioned Assignment
% LaTeX Template
% Version 1.0 (5/5/12)
%
% This template has been downloaded from:
% http://www.LaTeXTemplates.com
%
% Original author:
% Frits Wenneker (http://www.howtotex.com)
%
% License:
% CC BY-NC-SA 3.0 (http://creativecommons.org/licenses/by-nc-sa/3.0/)
%
%%%%%%%%%%%%%%%%%%%%%%%%%%%%%%%%%%%%%%%%%

%----------------------------------------------------------------------------------------
%	PACKAGES AND OTHER DOCUMENT CONFIGURATIONS
%----------------------------------------------------------------------------------------

\documentclass[paper=a4, fontsize=11pt]{scrartcl} % A4 paper and 11pt font size

\usepackage[T1]{fontenc} % Use 8-bit encoding that has 256 glyphs
\usepackage{fourier} % Use the Adobe Utopia font for the document - comment this line to return to the LaTeX default
\usepackage[english]{babel} % English language/hyphenation
\usepackage{amsmath,amsfonts,amsthm,amssymb} % Math packages
\usepackage{caption}
\usepackage{lipsum} % Used for inserting dummy 'Lorem ipsum' text into the template
\usepackage{float}
\restylefloat{table}

\usepackage{sectsty} % Allows customizing section commands
\sectionfont{\centering \normalfont} % Make all sections centered, the default font and small caps
\subsectionfont{\centering \normalfont} % Make all sections centered, the default font and small caps
\subsubsectionfont{\raggedright \normalfont}

\usepackage{fancyhdr} % Custom headers and footers
\pagestyle{fancyplain} % Makes all pages in the document conform to the custom headers and footers
\fancyhead{} % No page header - if you want one, create it in the same way as the footers below
\fancyfoot[L]{} % Empty left footer
\fancyfoot[C]{} % Empty center footer
\fancyfoot[R]{\thepage} % Page numbering for right footer
\renewcommand{\headrulewidth}{0pt} % Remove header underlines
\renewcommand{\footrulewidth}{0pt} % Remove footer underlines
\setlength{\headheight}{13.6pt} % Customize the height of the header
\newcommand*{\oldneg}{\mathord{\sim}}

\setlength\parindent{0pt} % Removes all indentation from paragraphs - comment this line for an assignment with lots of text

%----------------------------------------------------------------------------------------
%	TITLE SECTION
%----------------------------------------------------------------------------------------

\newcommand{\horrule}[1]{\rule{\linewidth}{#1}} % Create horizontal rule command with 1 argument of height

\title{	
\normalfont \normalsize 
\textsc{Discrete Math} \\ [25pt] % Your university, school and/or department name(s)
\horrule{0.5pt} \\[0.2cm] % Thin top horizontal rule
\large Homework 1 \\ % The assignment title
\horrule{2pt} \\[0.2cm] % Thick bottom horizontal rule
}
\author{Michael Introini} % Your name

\date{\normalsize{May 29, 2017}} % Today's date or a custom date

\begin{document}

\maketitle % Print the title

%----------------------------------------------------------------------------------------
%	SECTION 1.1
%----------------------------------------------------------------------------------------

\section*{Chapter 1}

\horrule{0.5pt} \\ % Thin top horizontal rule
\subsection*{Section 1.1}
\subsubsection*{Problem 12e}
\textbf{Let p, q, and r be the propositions}\\
\emph{p} : You have the flu.\\
\emph{q} : You miss the final examination.\\
\emph{r} : You pass the course.\\

Express the proposition as an English sentence

\textbf{e)} $(p \to \oldneg r) \lor (q \to \oldneg r)$\\

If you have the flue or you miss the final examination, you will \emph{not} pass the course.


\subsubsection*{Problem 14d,e}
\textbf{Let p, q, and r be the propositions}\\
\emph{p} : You get an A on the final exam.\\
\emph{q} : You do every exercise in this book.\\
\emph{r} : You get an A in this class.\\

Write these propositions using p, q, and r and logical connectives (including negations).\\
\textbf{d)} You get an A on the final, but you don't do every ex
ercise in this book; nevertheless, you get an A in this class.\\

$\oldneg q \land (p \land r)$\\

\textbf{e)} Getting an A on the final and doing every exercise in\\
this book is sufficient for getting an A in this class.\\

$(p \land q) \to r$\\

\subsubsection*{Problem 32d}
Construct a truth table for each of these compound propositions.\\
\textbf{d)} $(p \land q) \to (p \lor q)$\\

\begin{table}[H]
    \centering
    \begin{tabular}{ll|cc|c}
	$p$ & $q$ & \multicolumn{1}{l}{$(p \land q)$} & \multicolumn{1}{l|}{$(p \lor q)$} & \multicolumn{1}{l}{$(p \land q) \to (p \lor q)$} \\ \hline
	F & F & F & F & T \\
	F & T & F & T & T \\
	T & F & F & T & T \\
	T & T & T & T & T
    \end{tabular}
\end{table}

%----------------------------------------------------------------------------------------
%	SECTION 1.2
%----------------------------------------------------------------------------------------

\horrule{0.5pt} \\ % Thin top horizontal rule
\subsection*{Section 1.2}
\subsubsection*{Problem 36}

Four friends have been identified as suspects for an unauthorized access into a computer 
system. They have made statements to the investigating authorities. Alice said
\textit{Carlos did it}. John said \textit{I did not do it.} Carlos said
\textit{Diana did it.} Diana said \textit{Carlos lied when he said that
I did it.}\\

\textbf{a)} If the authorities also know that exactly one of the
four suspects is telling the truth, who did it? Explain
your reasoning.\\

John was the one that did it.\\
\textbf{Explination:}\\
Both Diana and Carlos cannot lie at the same time, which means one of them is
telling the truth. That implies that Alice is lying, and so is John. Since John is
the only one referring to himself, he is the one who did it.\\

\textbf{b)} If the authorities also know that exactly one is lying,
who did it? Explain your reasoning.\\

Carlos did it.\\
\textbf{Explination:}\\
Conversely, Carlos and Diana cannot tell the truth at the same time, and since only one
suspect is lying, it means that Alice is telling the truth which implies that Carlos did it.\\

%----------------------------------------------------------------------------------------
%	SECTION 1.3
%----------------------------------------------------------------------------------------

\horrule{0.5pt} \\ % Thin top horizontal rule
\subsection*{Section 1.3}
\subsubsection*{Problem 8b}

\textbf{8} Use De Morgan's laws to find the negation of each of the
following statements.

\textbf{b)} Yoshiko knows Java and calculus.\\

Yoshiko does not know Java, or does not know Calculus, or does not know both.

$\oldneg p \lor \oldneg q \lor \oldneg r$


\subsubsection*{Problem 10}
Show that each of these conditional statements is a tautology by using truth tables.

\textbf{b)} $[(p \to q) \land (q \to r)] \to (p \to r)$\\

\begin{table}[H]
\centering
\begin{tabular}{ccc|cccc|c}
    $p$ & $q$ & $r$ & $(p \to q)$ & $(q \to r)$ & $[(p \to q) \land (q \to r)]$ & $(p \to r)$ & $[(p \to q) \land (q \to r)] \to (p \to r)$ \\ \hline
    F & F & F & T & T & T & T & T \\
    F & F & T & T & T & T & T & T \\
    F & T & F & T & F & F & T & T \\
    F & T & T & T & T & T & T & T \\
    T & F & F & F & T & F & T & T \\
    T & F & T & F & T & F & T & T \\
    T & T & F & T & F & F & F & T \\
    T & T & T & T & T & T & T & T
\end{tabular}
\end{table}


\subsubsection*{Problem 32}
\textbf{32} Show that $(p \land q) \to r$ and $(p \to r) \land (q \to r)$ are not
logically equivalent.\\


\begin{table}[H]
    \centering
    \begin{tabular}{lll}
        $(p \land q) \to r$ &  & $(p \to r) \land (q \to r)$ \\
        $\oldneg (p \land q) \lor r$ &  & $(\oldneg p \lor r) \land (\oldneg q \lor r)$ \\
        $\oldneg p \lor \oldneg q \lor r$ & $\ne$ & $(\oldneg p \land \oldneg q) \lor r$
    \end{tabular}
\end{table}

\subsubsection*{Problem 54}
\textbf{54} Show that $p  |  (q  |  r)$ and $(p  |  q)  |  r$ are not equivalent,
so that the logical operator $|$ is not associative.

%--------------------------------------------------------------------------------------%
%	SECTION 1.4
%--------------------------------------------------------------------------------------%

\horrule{0.5pt} \\ % Thin top horizontal rule
\subsection*{Section 1.4}
\subsubsection*{Problem 14}

\textbf{14} Determine the truth value of each statement\\
\textbf{a)} $\exists x(x^3 = -1)$\\
\textbf{b)} $\exists x(x^4 < x^2)$\\
\textbf{c)} $\forall x((-x^2) = x^2)$\\
\textbf{d)} $\forall x(2x > x)$\\

\textbf{a)} True\\
\textbf{b)} True\\
\textbf{c)} True\\
\textbf{d)} False\\

\subsubsection*{Problem 36}
Find a counterexample, if possible, to these universally
quantified statements, where the domain for all variables
consists of all real numbers.\\

\textbf{a)} $\forall x(x^2 \ne x)$\\
\textbf{b)} $\forall x(x^2 \ne 2)$\\
\textbf{c)} $\forall x(|x| > 0)$\\

\textbf{a)} $\exists x(x^2 = x)$; Because when $x=0$, $0^2=0$.\\ 
Therefor $x=0$ is a counterexample\\
\textbf{a)} $\exists x(x^2 = 2)$; Because when $x = \sqrt{2}$, $(\sqrt{2})^2 = 2$.\\ 
Therefore $x=\sqrt{2}$ is a counterexample\\
\textbf{c)} $\exists x(|x| \le 0)$; Because when $x=0$, $|0| \le 0$ is True.\\
Therefore $x=0$ is a counterexample.

%--------------------------------------------------------------------------------------%
%	SECTION 1.5
%--------------------------------------------------------------------------------------%

\horrule{0.5pt} \\ % Thin top horizontal rule
\subsection*{Section 1.5}
\subsubsection*{Problem 4}

Let $P(x, y)$ be the statement “Student x has taken class
y,” where the domain for $x$ consists of all students in your
class and for $y$ consists of all computer science courses
at your school. Express each of these quantifications in
English.\\

\textbf{b)} $\exists x \forall y P(x, y)$\\
\textbf{c)} $\forall x \exists y P(x, y)$\\

\textbf{b)} There is a student $x$ who has taken every Computer Science class in my school.\\
\textbf{c)} For every student $x$ this exists at least one Computer Science class in my school.

\subsubsection*{Problem 24}

Translate each of these nested quantifications into an English 
statement that expresses a mathematical fact. The
domain in each case consists of all real numbers.\\

\textbf{a)} $\exists x \forall y(x + y = y)$\\
\textbf{b)} $\forall x \forall y((x \ne 0) \land (y \ne 0) \leftrightarrow (xy \ne 0))$\\

\textbf{a)} There exists at least one real number $x$ such that the sum of its value and a real number $y$ equal $y$.\\
\textbf{b)} For every pair $x,y$ their product is non-zero if and only if $x$ and $y$ are also non-zero.

\subsubsection*{Problem 28}

Determine the truth value of each of these statements if
the domain of each variable consists of all real numbers.

\textbf{a)} $\forall x \exists y(x^2 = y)$\\
\textbf{b)} $\forall x \exists y(x = y^2)$\\

\textbf{a)} True\\
\textbf{b)} False

\subsubsection*{Problem 30}
Rewrite each of these statements so that negations appear 
only within predicates (that is, so that no negation
is outside a quantifier or an expression involving logical
connectives).

\textbf{b)} $\oldneg \forall x \exists y P(x, y)$\\
\textbf{c)} $\oldneg \exists y(Q(y) \land \forall x \oldneg R(x, y))$


\textbf{b)} $\exists x \forall y \oldneg P(x, y)$\\
\textbf{c)} $\forall y(\oldneg Q(y) \lor \exists x R(x, y))$\\

%--------------------------------------------------------------------------------------%
%	SECTION 1.6
%--------------------------------------------------------------------------------------%

\horrule{0.5pt} \\ % Thin top horizontal rule
\subsection*{Section 1.6}
\subsubsection*{Problem 10}
For each of these sets of premises, what relevant conclusion 
or conclusions can be drawn? Explain the rules of inference 
used to obtain each conclusion from the premises.\\

\textbf{c)} 
"All insects have six legs." "Dragonflies are insects."
"Spiders do not have six legs." "Spiders eat dragonflies."\\

All insects have six legs.\\
Dragonflies are insects.\\
Spiders do not have six legs.\\
Spiders eat dragonflies.\\

\textbf{Hypothetical Syllogism:} Dragonflies have six legs.\\
\textbf{Modus Tollens:} Spiders are not insects.\\
\textbf{Hypothetical Syllogism:} Spiders eat all insects.\\

\textbf{Conclusion:} Spiders eat insects and dragonflies have six legs.

\subsubsection*{Problem 12}
Show that the argument form with premises 
$(p \land t) \to (r \lor s)$, $q \to (u \land t)$, $u \to p$, and $\oldneg s$ and conclusion
$q \to r$ is valid by first using Exercise 11 and then using rules of inference from Table 1.\\

$r$ Modus Ponens

%--------------------------------------------------------------------------------------%
%	SECTION 1.7
%--------------------------------------------------------------------------------------%

\horrule{0.5pt} \\ % Thin top horizontal rule
\subsection*{Section 1.7}
\subsubsection*{Problem 8}
\textbf{8)} Prove that if $n$ is a perfect square, then $n + 2$ is not a
perfect square.

\textbf{Contradiction}
$p = n$ and $q = n+2$\\
Assume $n = a^2$, and $n+2 = b^2$\\
$2 = b^2 - n$\\
$2 = b^2 - a^2$ (Replacing $n$ by $a^2$)\\
$2 = (a+b)(a-b)$\\
$2 = (a+b)$ and $1=(a-b)$, which added together is $2b = 3$\\
$b = \frac{3}{2}$ which is not a perfect square.\\
\\
\textbf{c)} Prove or disprove that the product of a nonzero rational
number and an irrational number is irrational.\\

\textbf{Contradiction}\\
Let $p = \frac{a}{b}$ where $a \not = 0$, and $b \not = 0$, and $k$ be an arbitrary irrational number.\\
$\frac{a}{b} * k = \frac{ak}{b}$\\

%--------------------------------------------------------------------------------------%
%	SECTION 1.8
%--------------------------------------------------------------------------------------%

\horrule{0.5pt} \\ % Thin top horizontal rule
\subsection*{Section 1.8}
\subsubsection*{Problem 12}
Show that the product of two of the numbers $65^{1000} - 8^{2001} + 3^{177}$ , 
$79^{1212} - 9^{2399} + 2^{2001}$ , and $24^{4493} - 5^{8192} + 7^{1777}$ 
is nonnegative. Is your proof constructive or nonconstructive? 
[Hint: Do not try to evaluate these numbers!]\\

Let $a, b,$ and $c$ represent the three summations above.
Since either $a, b,$ or $c$ can be a negative, positive, or zero value number, we have
a few cases to consider.\\
If either of the two numbers is zero, then our product is zero; which is nonnegative.\\
If we chose 2 negative numbers, our product is a positive number.\\
For the case where we have a positive and negative number, we may end up with either a negative or
positive number depending on our choice. Therefore, as long as we choose the right value, we can
ensure that e always get a nonnegative value.\\

Non-Constructive\\
\subsubsection*{Problem 36}
Prove that between every rational number and every irrational number there is an irrational number.\\

Let $r = \frac{a}{b}$ where $b \not = 0$, also let $\sqrt{2}r$ be an arbitrary irrational number.\\
Since know that $r < \sqrt{2}r$, we can infer that $r < \frac{1}{2}\sqrt{2}r < \sqrt{2}r$\\
This proves that between every rational and irrational number, there is another irrational number.\\

%--------------------------------------------------------------------------------------%
%	SECTION 2.1
%--------------------------------------------------------------------------------------%

\horrule{0.5pt} \\ % Thin top horizontal rule
\subsection*{Section 2.1}
\subsubsection*{Problem 18}
Find two sets $A$ and $B$ such that $A \in B$ and $A \subseteq B$

Let $A = \{0,1,2,3\}$ and $B = \{0,1,2,3,4,5,6,7\}$\\
\subsubsection*{Problem 20}
What is the cardinality of each of these sets?\\
\textbf{d)} $\{\emptyset , \{\emptyset \}, \{\emptyset \}, \{\emptyset \}\}\}$\\
3
%--------------------------------------------------------------------------------------%
%	SECTION 2.2
%--------------------------------------------------------------------------------------%

\horrule{0.5pt} \\ % Thin top horizontal rule
\subsection*{Section 2.2}
\subsubsection*{Problem 18}
Let A, B, and C be sets. Show that\\
\textbf{d)} $(A - C) \cap (C - B) = \emptyset$\\ 
$(A - C) \cap (C - B) = \emptyset$\\
$(A \cap \overline{C}) \cap (C \cap \overline{B})$\\

\subsubsection*{Problem 32}
Find the symmetric difference of $\{1, 3, 5\}$ and $\{1, 2, 3\}$\\
$\{2,5\}$\\

\subsubsection*{Problem 48}
Let $A_i = \{\dots,-2,-1,0,1,\ldots,i\}$. Find\\
\textbf{a)} $\displaystyle \bigcup_{i=1}^{n}A_i$ 
\textbf{b)} $\displaystyle \bigcap_{i=1}^{n}A_i$\\

\textbf{a)} $\displaystyle \bigcup_{i=1}^{n}A_1 = \{\ldots,-2,-1,0,1,\ldots n\}$\\
\textbf{b)} $\displaystyle \bigcap_{i=1}^{n}A_1 = \{\ldots,-2,-1,0,1\}$\\ 

\end{document} 
